Пусть \(G=(V_1,V_2,E)\) — двудольный граф с долями \(V_1,V_2\).
\begin{lemma}[Слабая двойственность]
В любом графе \(\alpha'(G)\le\beta(G)\).
\end{lemma}
\begin{proof}
Пусть \(M\) — паросочетание, \(C\) — вершинное покрытие. Каждое ребро \(M\)
имеет конец в \(C\), а одна вершина \(C\) покрывает не более одного ребра \(M\)
(рёбра \(M\) не имеют общих концов). Значит, \(|M|\le|C|\).
\end{proof}
Для треугольника \(\alpha'(K_3)=1<2=\beta(K_3)\), так что в общем случае
равенства нет. В двудольном графе оно есть.

\begin{theorem}[Кёниг, 1931]
Если \(G\) двудолен, то \(\alpha'(G)=\beta(G)\): размер максимального
паросочетания равен размеру минимального вершинного покрытия.
\end{theorem}
\begin{proof}
По лемме достаточно построить вершинное покрытие \(B\) размера \(|M|\),
где \(M\) — максимальное паросочетание.

Пусть \(U_1\subseteq V_1\) — свободные вершины доли \(V_1\), \(U_2\subseteq V_2\) —
свободные вершины доли \(V_2\). Назовём вершину \emph{достижимой}, если до неё
есть \(M\)-чередующийся путь из \(U_1\), начинающийся ребром не из \(M\)
(вершины \(U_1\) достижимы путём длины 0). Положим
\[
Y_1=\text{достижимые из }V_1\setminus U_1,\quad Z_1=V_1\setminus(U_1\cup Y_1),\qquad
Y_2=\text{достижимые из }V_2,\quad Z_2=V_2\setminus Y_2 .
\]
\textit{Свойства.}
\begin{enumerate}\tightlist
\item \(Y_2\cap U_2=\varnothing\): иначе путь до свободной \(y\in Y_2\) был бы
  \(M\)-дополняющим, что невозможно по теореме Бержа. Значит, \(U_2\subseteq Z_2\).
\item Путь из \(U_1\) переходит из \(V_1\) в \(V_2\) по рёбрам не из \(M\), а из
  \(V_2\) в \(V_1\) — по рёбрам \(M\). Поэтому партнёр по \(M\) каждой
  \(y\in Y_2\) лежит в \(Y_1\) (путь продолжается ребром \(M\)), и каждая
  вершина \(Y_1\) получена так. Итак, \(M\) — биекция \(Y_1\leftrightarrow Y_2\);
  следовательно, партнёры вершин \(Z_1\) (они покрыты по определению) лежат в \(Z_2\).
\item Если \(x\in U_1\cup Y_1\) и \(xy\in E\setminus M\), то \(y\in Y_2\): путь до \(x\)
  кончается ребром \(M\) (или пуст), и его можно продлить ребром \(xy\).
\end{enumerate}
Положим \(B=Z_1\cup Y_2\).

\emph{\(B\) — покрытие.} Ребро \(xy\) (\(x\in V_1\), \(y\in V_2\)) с \(x\in Z_1\)
покрыто. Если \(x\in U_1\cup Y_1\): при \(xy\notin M\) имеем \(y\in Y_2\) по
свойству 3; при \(xy\in M\) вершина \(x\in Y_1\) и её партнёр \(y\in Y_2\) по
свойству 2.

\emph{\(|B|=|M|\).} Каждое ребро \(M\) соединяет либо \(Y_1\) с \(Y_2\), либо
\(Z_1\) с \(Z_2\) (свойство 2), т.\,е. имеет ровно один конец в \(B\);
и каждая вершина \(B\) покрыта (свойство 1 и определение \(Z_1\)).
Значит, \(|B|=|M|=\alpha'(G)\), и \(\beta(G)\le\alpha'(G)\).
\end{proof}
\begin{graphillustration}{Структура из доказательства Кёнига. Толстые рёбра — $M$, крестиком отмечены невозможные рёбра. Покрытие $B=Z_1\cup Y_2$ выделено двойным контуром групп.}{fig:02-konig}
\begin{tikzpicture}[font=\small]
\node[gi free] (u1) at (0,2) {};
\node[gi vertex] (y1a) at (2.2,2) {};
\node[gi vertex] (y1b) at (3.2,2) {};
\node[gi vertex] (z1a) at (5.6,2) {};
\node[gi vertex] (z1b) at (6.6,2) {};
\node[gi free] (u2) at (0,0) {};
\node[gi vertex] (y2a) at (2.2,0) {};
\node[gi vertex] (y2b) at (3.2,0) {};
\node[gi vertex] (z2a) at (5.6,0) {};
\node[gi vertex] (z2b) at (6.6,0) {};
\draw[gi match] (y1a) -- (y2a); \draw[gi match] (y1b) -- (y2b);
\draw[gi match] (z1a) -- (z2a); \draw[gi match] (z1b) -- (z2b);
\draw[gi edge] (u1) -- (y2a); \draw[gi edge] (y1a) -- (y2b);
\draw[gi edge] (z1a) -- (y2b); \draw[gi edge] (z1b) -- (z2a);
\draw[gi removed] (u1) -- (u2) node[pos=.5,fill=white,inner sep=1pt] {$\times$};
\draw[gi removed] (y1b) -- (z2a) node[pos=.5,fill=white,inner sep=1pt] {$\times$};
\node[gi group,fit=(u1),label=above:$U_1$] {};
\node[gi group,fit=(y1a)(y1b),label=above:$Y_1$] {};
\node[gi group,double,fit=(z1a)(z1b),label=above:{$Z_1\subseteq B$}] {};
\node[gi group,fit=(u2),label=below:$U_2$] {};
\node[gi group,double,fit=(y2a)(y2b),label=below:{$Y_2\subseteq B$}] {};
\node[gi group,fit=(z2a)(z2b),label=below:$Z_2$] {};
\node[gi note,anchor=west] at (7.6,2) {$V_1$};
\node[gi note,anchor=west] at (7.6,0) {$V_2$};
\node[gi note,anchor=west] at (8.4,1) {$U_1$—$U_2$: дополняющий путь;\\$U_1\cup Y_1$—$Z_2$: $Z_2$ стала бы\\достижимой.};
\end{tikzpicture}
\end{graphillustration}


\begin{cor}[Двудольный граф, \(\delta>0\)]
\(\alpha(G)=\beta'(G)\).
\end{cor}
\begin{proof}
По билету~\ref{ticket:3} \(\alpha+\beta=n=\alpha'+\beta'\), а \(\beta=\alpha'\) по Кёнигу.
\end{proof}
\begin{example}
В \(C_6\) (двудолен) \(\alpha'=3\) и покрытие из трёх вершин одной доли.
В \(C_5\) (не двудолен) \(\alpha'=2\), а \(\beta=3\): две вершины покрывают
не более четырёх рёбер из пяти.
\end{example}
